\subsection{An abstract cyclic norm system with unbounded growth}

\begin{proposition}[KT-ORDER-FIVE-ABSTRACT-NORM-LIMIT]
Let $p\equiv1\pmod {10}$ be prime and identify
$\Lambda=\mathbf Z_p[[\Gamma]]=\mathbf Z_p[[T]]$ for
$\Gamma\simeq\mathbf Z_p$ with a topological generator $\gamma=1+T$.
There is an explicit continuous $\Lambda$-module $X$, with a commuting
action of $C_5$ through one primitive Teichm\"uller character, for which
the finite quotients $A_n=X/\omega_n(T)X$, where
$\omega_n(T)=(1+T)^{p^n}-1$, have the following properties:
\[
 A_n\simeq\mathbf Z/p^{n+1}\mathbf Z,\qquad
 |A_n^{\Gamma}|=p,\qquad
 j_{0,n}:A_0\lhook\joinrel\longrightarrow A_n,
 \qquad j_{0,n}(A_0)=A_n^{\Gamma}.                     \tag{O5M.1}
\]
The natural norm and transfer maps satisfy their degree and algebraic
norm identities.  Nevertheless $|A_n|$ is unbounded, and $X$ has
$\mu=0$, $\lambda=1$.
This is an abstract module example, not a class group of a number field.
\end{proposition}

\begin{proof}
Take
\[
             X=\Lambda/(T+p)\simeq\mathbf Z_p,
             \qquad \gamma=1+T\ \text{acts as }1-p.
                                                               \tag{O5M.2}
\]
For odd $p$, the lifting-the-exponent identity gives
\[
       v_p\bigl((1-p)^{p^n}-1\bigr)=n+1.               \tag{O5M.3}
\]
For clarity, this also follows inductively: if $v_p(z)=r\ge1$,
then the term $pz$ in $(1+z)^p-1$ has valuation $r+1$;
all middle terms have valuation at least $1+2r>r+1$, and
$z^p$ has valuation $pr>r+1$.  Start with $z=-p$.
Thus $A_n$ has the order and abstract structure in (O5M.1).
The action of $\gamma$ on it factors through a group of order $p^n$.

For $m\ge n$, put
\[
 \nu_{m,n}(T)=\frac{\omega_m(T)}{\omega_n(T)}
       =\sum_{i=0}^{p^{m-n}-1}(1+T)^{ip^n}.
                                                               \tag{O5M.4}
\]
Define $N_{m,n}:A_m\to A_n$ by the natural quotient and
$j_{n,m}:A_n\to A_m$ by multiplication by $\nu_{m,n}(-p)$.
These are well-defined because $\omega_n$ divides $\omega_m$.
Equation (O5M.3) gives
\[
              v_p(\nu_{m,n}(-p))=m-n.                 \tag{O5M.5}
\]
Every $N_{m,n}$ is onto and every $j_{n,m}$ is injective.
The maps compose correctly because the quotients $\nu_{m,n}$
multiply along a tower.  Moreover,
\[
 N_{m,n}j_{n,m}=p^{m-n}\quad\text{on }A_n,
 \qquad
 j_{n,m}N_{m,n}=\sum_{i=0}^{p^{m-n}-1}\gamma^{ip^n}
       \quad\text{on }A_m.                            \tag{O5M.6}
\]
For the first identity, $(1+T)^{p^n}=1$ on $A_n$;
the second is precisely the defining polynomial (O5M.4).

Since $\gamma-1$ acts as $-p$, its kernel on $A_n$ is the
subgroup of order $p$.  By (O5M.5), this is exactly the image
of $A_0$ under $j_{0,n}$, proving (O5M.1).
The inverse limit under the quotient norms is $\mathbf Z_p=X$.
Its distinguished characteristic polynomial is $T+p$, giving
$\mu=0$ and $\lambda=1$.  In the direct system of transfers,
the nonzero class of $A_0$ survives because every transfer is
injective.

Finally, choose a primitive fifth root of unity
$\xi\in\mathbf Z_p^{\times}$, which exists because $5\mid p-1$,
and let a generator of $C_5$ act by multiplication by $\xi$.
This commutes with $T$ and all the displayed maps.  Each $A_n$
therefore has just that one primitive order-five character.
\end{proof}

\paragraph{Exact limitation.}
The base component has length one, the ambiguous subgroup has the
same order at every level, and all the stated norm/transfer identities
hold, but bounded growth fails.  Thus those module data alone do not
force the extra unit-norm condition in
KT-ORDER-FIVE-GREENBERG-UNIT-GATE.
No realization by class fields, global units, or an actual
cyclotomic number-field tower is asserted.  In particular, this
example is not a counterexample to Greenberg's conjecture and does
not remove an arithmetic class component.
